\documentclass[11pt,a4paper]{article}
\usepackage{rclattice-report}

% paths resolve from doc/, which is where build.sh runs pdflatex
\graphicspath{{../output/wall/}{reports/kutay_wall/figures/}}


\title{\vspace{-14mm}\textbf{Kutay's Wall}\\[2mm]
\large Lattice Model, Cyclic Analysis Parameters and Results\\[1mm]
\normalsize Reversed-cyclic analysis to 4.0\% drift ratio}
\author{\normalsize \texttt{rclattice} --- \texttt{examples/wall/}}
\date{\normalsize \today}

\begin{document}
\maketitle
\thispagestyle{empty}

\section*{Specimen}

Specimen SW-NC-FF from Şahinkaya, Binbir, Orakcal \& Ilki (2025), \emph{Behavior of Nonconforming
Flexure-Controlled RC Structural Walls Under Reversed Cyclic Lateral Loading}, \emph{Buildings}
\textbf{15}, 4501. Full-scale nonconforming flexure-controlled structural wall with low-strength
concrete, plain longitudinal reinforcement and no boundary confinement.

Unit system throughout: \textbf{N, mm, MPa, s}.

\section{Geometry}

Every dimension is taken directly from the paper and used unchanged; none is derived, so the
modelled value equals the reported one in every row below.

\begin{tabularx}{\textwidth}{@{}p{46mm}p{34mm}>{\raggedright\arraybackslash}X@{}}
\toprule
\textbf{Quantity} & \textbf{Value} & \textbf{Source in the paper} \\
\midrule
Wall length $l_w$ (in-plane)        & 1000\unit{mm} & Sec.~2.1 \\
Wall height $h_w$                   & 3000\unit{mm} & Sec.~2.1; $h_w/l_w = 3.0$ \\
Wall thickness $t_w$ (out-of-plane) & 200\unit{mm}  & Sec.~2.1 \\
Shear span $a$                      & 2200\unit{mm} & Sec.~2.2, ``lateral loading was applied 2.2\unit{m} above the pedestal''; $a/l_w = 2.2$, matching the reported SSDR \\
Pedestal, in-plane                  & $1900 \times 600$\unit{mm} & Sec.~2.1, length $\times$ depth \\
Pedestal, out-of-plane              & 800\unit{mm} & Sec.~2.1, width; carried as an equivalent modulus (\S\ref{sec:ideal}) \\
Upper section (loading head)        & $y = 2000$--$3000$\unit{mm} & Sec.~2.1, ``the top 1\unit{m}'', reduced bar spacing and a higher-strength cast \\
Model origin                        & $x=0$ at wall mid-length, $y=0$ at the W--P interface & model convention, not from the paper \\
\bottomrule
\end{tabularx}

The domain is a compound of two axis-aligned rectangles --- pedestal $(-950,-600)$ to $(950,0)$ and
wall $(-500,0)$ to $(500,3000)$ --- meshed separately and merged at coincident nodes, so the
wall--pedestal interface is continuous.

\section{Discretization}

Discretization has no counterpart in the paper; every entry below is a modelling choice. The one
constraint the specimen imposes is that the grid must place a node on every bar line, which
restricts $d$ to divisors of the 200\unit{mm} bar spacing and of the 450\unit{mm} boundary-bar
offset.

\begin{tabularx}{\textwidth}{@{}p{50mm}p{38mm}>{\raggedright\arraybackslash}X@{}}
\toprule
\textbf{Parameter} & \textbf{Value} & \textbf{Note} \\
\midrule
Mesh size $d$          & 50\unit{mm}   & divides every reinforcement coordinate \\
Node generation        & gmsh, structured transfinite grid & one node source for the whole domain \\
Horizon $\delta$       & $1.5\,d = 75$\unit{mm} & 8 neighbours per interior node \\
Strut connectivity     & every node pair with separation $\le \delta$, deduplicated & orthogonal ($d$) + diagonal ($d\sqrt2$) \\
Element type           & \texttt{Truss} (2-node, translational DOF only) & \\
Degrees of freedom     & $ndm = ndf = 2$ & \\
Nodes                  & 1767 & \\
Concrete struts        & 6736 & \\
Reinforcement struts   & 896  & \\
Uniaxial materials     & 10   & concrete deduplicated by (zone, strut length) \\
\bottomrule
\end{tabularx}

Strut lengths present on the regular grid are $d = 50$\unit{mm} (orthogonal) and
$d\sqrt2 = 70.71$\unit{mm} (diagonal).

\section{Strut-area calibration}

Strut areas are set by the overlapping-lattice energy balance of Aydın (2017), \emph{Overlapping
lattice modeling for concrete fracture simulations using sequentially linear analysis}, MSc thesis,
METU, Section 2.2, Equations 2.1--2.3. An affine strain field is imposed, the continuum's stored
elastic energy is equated to the lattice's evaluated with $EA = 1$, and the uniform axial rigidity
follows as $EA_t = W_{\text{continuum}} / W_{\text{lattice}}$. The affine field makes every strut
elongation exact, so no finite element solve and no reference model is involved.

\begin{tabularx}{\textwidth}{@{}p{50mm}p{38mm}>{\raggedright\arraybackslash}X@{}}
\toprule
\textbf{Quantity} & \textbf{Value} & \textbf{Note} \\
\midrule
Calibrated strut area $A_t$      & 6042.9\unit{mm$^2$} & uniform, all concrete struts \\
Normalised                       & $0.6043 \times (t_w \cdot d)$ & $t_w \cdot d = 10\,000$\unit{mm$^2$} \\
Poisson ratio used in Eq. 2.1    & 0.20 & assumed; the paper reports $E_c$ per ASTM C469 but no measured $\nu$ \\
Lattice's self-consistent $\nu$  & 0.1796 & where the normal-strain and pure-shear balances agree \\
Shear-stiffness error at $\nu=0.20$ & 2.49\% & \\
Grid anisotropy $|A_x - A_y|/A_x$ & 1.90\% & \\
Calibration patch                & wall rectangle $1000 \times 3000$\unit{mm} & pedestal excluded \\
\bottomrule
\end{tabularx}

$A_t$ is independent of $E$, so the single value applies to all three concrete zones; each zone's
modulus is carried by its material. $A_t/(t_w d)$ is a constant of the ($\delta$, $\nu$) lattice.

\section{Material models}

\subsection{Concrete --- \texttt{Concrete02}, crack-band regularized}

Concrete struts use \texttt{Concrete02} with both softening branches sized from the strut length
$L$, holding the fracture energy per unit crack area constant:
\[
E_{ts} = \frac{f_t^2 L}{2 G_f} \ \ (\text{capped at } 0.5E), \qquad
\varepsilon_U = \varepsilon_{c0} + \frac{2 G_{fc}}{(f_c + f_{cu})L}, \qquad
G_{fc} = 250\,G_f .
\]
The residual crushing strength is floored at $0.2 f_c$. Tensile fracture energies follow CEB-FIP
MC90, $G_f = 0.030 (f_{cm}/10)^{0.7}$ with $f_{cm} = f_c + 8$\unit{MPa} and 16\unit{mm} maximum
aggregate.

\begin{center}\small
\begin{tabular}{@{}lrrrrrr@{}}
\toprule
\textbf{Zone} & $E$ & $f_c$ & $\varepsilon_{c0}$ & $f_{cu}$ & $f_t$ & $G_f$ \\
 & (MPa) & (MPa) & (--) & (MPa) & (MPa) & (N/mm) \\
\midrule
Test unit ($0 \le y < 2000$)   & 16\,100 & 15.1 & 0.0022 & 4.0  & 1.5 & 0.053 \\
Upper section ($y \ge 2000$)   & 26\,400 & 43.4 & 0.0025 & 10.0 & 3.3 & 0.093 \\
Pedestal ($y < 0$)             & 93\,600 & \multicolumn{5}{l}{linear \texttt{Elastic}; see \S\ref{sec:ideal}} \\
\bottomrule
\end{tabular}\end{center}

Provenance of each column:

\begin{tabularx}{\textwidth}{@{}p{22mm}p{40mm}>{\raggedright\arraybackslash}X@{}}
\toprule
\textbf{Quantity} & \textbf{Origin} & \textbf{Value in the paper} \\
\midrule
$f_c$ & measured, used unchanged & Sec.~2.1: 15.1\unit{MPa} test unit, 43.4\unit{MPa} upper section,
33.7\unit{MPa} pedestal --- cylinder averages per ASTM C39, curves in Fig.~6 \\
$E$   & measured, used unchanged & Sec.~2.1, Eq.~(1) (secant between $\varepsilon_1 = 0.00005$ and
40\% of ultimate load): 16.1\unit{GPa} test unit, 26.4\unit{GPa} upper section, 23.4\unit{GPa}
pedestal. The 93\,600\unit{MPa} pedestal value is $23\,400 \times 800/200$ (\S\ref{sec:ideal}) \\
$\varepsilon_{c0}$ & read from the measured curves & Fig.~6b (test unit) and 6a (upper section);
strain at peak stress \\
$f_{cu}$ & residual plateau, floored at $0.2f_c$ & not tabulated in the paper; Fig.~6b shows the
test-unit curve falling to about 5--6\unit{MPa} by $\varepsilon_c = 0.010$ \\
$f_t$ & assumed; not measured & no tension test is reported. The upper-section and pedestal values
(3.3 and 2.7\unit{MPa}) match MC90 $f_{ctm} = 1.4((f_c-8)/10)^{2/3}$, which gives 3.25 and 2.63.
The test unit's 1.5\unit{MPa} is above that formula's 1.11\unit{MPa} \\
$G_f$ & derived; not measured & CEB-FIP MC90, $G_f = 0.030(f_{cm}/10)^{0.7}$ with
$f_{cm} = f_c + 8$\unit{MPa} and 16\unit{mm} maximum aggregate. Aggregate size is not reported \\
\bottomrule
\end{tabularx}

Emitted \texttt{Concrete02} arguments $(f_{pc}, \varepsilon_{c0}, f_{pcu}, \varepsilon_U, \lambda,
f_t, E_{ts})$, per zone and strut length (compression negative):

\begin{center}\small
\begin{tabular}{@{}llrrrrrrr@{}}
\toprule
\textbf{Zone} & $L$ (mm) & $f_{pc}$ & $\varepsilon_{c0}$ & $f_{pcu}$ & $\varepsilon_U$ & $\lambda$ & $f_t$ & $E_{ts}$ \\
\midrule
Test unit & 50.00 & $-15.10$ & $-0.00220$ & $-4.00$  & $-0.02995$ & 0.1 & 1.50 & 1061.3 \\
Test unit & 70.71 & $-15.10$ & $-0.00220$ & $-4.00$  & $-0.02182$ & 0.1 & 1.50 & 1500.9 \\
Upper     & 50.00 & $-43.40$ & $-0.00250$ & $-10.00$ & $-0.01992$ & 0.1 & 3.30 & 2927.4 \\
Upper     & 70.71 & $-43.40$ & $-0.00250$ & $-10.00$ & $-0.01481$ & 0.1 & 3.30 & 4140.0 \\
\bottomrule
\end{tabular}\end{center}

\subsection{Reinforcement --- \texttt{Steel02}}

Arguments $(f_y, E_0, b, R_0, c_{R1}, c_{R2})$ with $R_0 = 18.0$, $c_{R1} = 0.925$,
$c_{R2} = 0.15$ (the \texttt{Steel02} defaults; the paper reports no cyclic bar test to fit them
against).

$f_y$, $E_0$, $f_{max}$ and $\varepsilon_{max}$ are the measured values of the paper's Table~2,
used unchanged. The strain-hardening ratio $b$ is not tabulated there and is computed from those
same measurements as the secant from yield to the maximum-stress point,
$b = (f_{max} - f_y)\,/\,[E_s(\varepsilon_{max} - \varepsilon_y)]$, using the reported
$\varepsilon_y$ (0.148\% for $\phi14$, 0.135\% for $\phi12$, 0.157\% for $\phi10$).

\begin{center}\small
\begin{tabular}{@{}lrrrrrl@{}}
\toprule
\textbf{Bar} & $f_y$ (MPa) & $E_0$ (MPa) & $b$ & $f_{max}$ (MPa) & $\varepsilon_{max}$ & \textbf{Use} \\
\midrule
$\phi14$ & 294 & 198\,600 & 0.0024 & 396 & 21.4\% & boundary longitudinal \\
$\phi12$ & 311 & 230\,300 & 0.0021 & 414 & 21.0\% & web vertical \\
$\phi10$ & 340 & 216\,600 & 0.0019 & 425 & 21.3\% & web horizontal \\
$\phi16$ & 294 & 198\,600 & 0.0024 & --- & --- & loading head ($\phi$14 properties used) \\
\bottomrule
\end{tabular}\end{center}

Table~2 covers $\phi8$, $\phi10$, $\phi12$ and $\phi14$. The $\phi16$ loading-head bars are not
tested there, so they carry the $\phi14$ properties. $\phi8$ (boundary hoops in the other specimen,
SW-IC-FF) is not used in this model.

\section{Reinforcement layout}

Two reinforcement curtains (front and back) are collapsed onto one in-plane line, so each modelled
bar line carries both curtains' area. Reinforcement shares lattice nodes with the concrete
(perfect bond).

\begin{tabularx}{\textwidth}{@{}p{30mm}p{36mm}r>{\raggedright\arraybackslash}X@{}}
\toprule
\textbf{Family} & \textbf{Position} & \textbf{Area (mm$^2$)} & \textbf{Detail and source} \\
\midrule
Boundary longitudinal & $x = \pm450$\unit{mm} & 615.75 & 4$\phi$14 (Fig.~3b, Fig.~4 Section 2-2), lumped at the boundary-region centroid \\
Web vertical          & $x = \pm100, \pm300$\unit{mm} & 226.19 & $\phi$12@200 (Fig.~3b), two curtains \\
Web horizontal        & $y = 0, 200, \dots, 3000$\unit{mm} & 157.08 & $\phi$10@200 (Fig.~3b), two curtains, full width \\
Loading head          & $y = 2500, \dots, 3000$\unit{mm} & 402.12 & $\phi$16@100 (Fig.~3c) \\
Boundary confinement  & --- & --- & none: Table~1 gives $\rho_{con} =$ ``NC'' and Sec.~2.1 states SW-NC-FF had no boundary transverse reinforcement \\
\bottomrule
\end{tabularx}

The boundary region is 10\% of $l_w$ (Table~1, boundary length ratio), i.e.\ 100\unit{mm}, so the
4$\phi$14 group centroid sits at $x = \pm450$\unit{mm}. Vertical bars are anchored into the
pedestal: each runs down to $y = -500$\unit{mm} and turns $300$\unit{mm} outward, away from the
wall centreline, following the hooked anchorage of Fig.~3b--c.

Reinforcement ratios, model against the paper's Table~1:

\begin{tabularx}{\textwidth}{@{}p{30mm}p{40mm}p{20mm}>{\raggedright\arraybackslash}X@{}}
\toprule
\textbf{Ratio} & \textbf{Model} & \textbf{Table~1} & \textbf{Note} \\
\midrule
$\rho_v$ (web vertical)   & $2A_{\phi12}/(t_w s) = 0.566\%$ & 0.57\% & agrees \\
$\rho_t$ (web horizontal) & $2A_{\phi10}/(t_w s) = 0.393\%$ & 0.39\% & agrees \\
$\rho_b$ (boundary)       & $615.75/(0.10 l_w t_w) = 3.08\%$ & 2.37\% & \textbf{unresolved discrepancy}
--- see below \\
\bottomrule
\end{tabularx}

The modelled $\rho_b$ does not reproduce Table~1. Taking 4$\phi$14 over a boundary region of
$0.10 l_w$ gives $615.75 / (100 \times 200) = 3.08\%$, against the reported 2.37\%. The same
calculation applied to the companion specimen SW-IC-FF --- 4$\phi$14 over $0.20 l_w$ --- gives
1.539\% against its reported 1.54\%, so the definition itself is consistent. Reproducing 2.37\%
would require about 474\unit{mm$^2$} (roughly 3$\phi$14) or a boundary length near 130\unit{mm}.
The model uses the 4$\phi$14 of Fig.~3b and Fig.~4; the resulting $+30\%$ of boundary steel acts
directly on flexural strength.

\section{Boundary conditions and loading}

\begin{tabularx}{\textwidth}{@{}p{36mm}p{28mm}>{\raggedright\arraybackslash}X@{}}
\toprule
\textbf{Item} & \textbf{Source} & \textbf{Definition} \\
\midrule
Support        & Sec.~2.2 & the pedestal was bolted to the strong floor with fourteen high-strength
rods and its sliding and rigid-body rotation measured as negligible; modelled as all nodes on the
pedestal soffit $y = -600$\unit{mm} fixed in both translations \\
Axial load     & Sec.~2.2, Eq.~(2) & $N = 20\%(t_w \times l_w \times f_c') \approx 600$\unit{kN}
compression, held constant; distributed over the wall top edge $y = 3000$\unit{mm} \\
Axial load ratio & Table~3 (ALR) & 20\%, reproduced exactly \\
Out-of-plane restraint & Sec.~2.2 & roller frames prevented out-of-plane displacement; the model is
2D in-plane, so this is satisfied by construction \\
Lateral drive  & Sec.~2.2 & applied 2.2\unit{m} above the pedestal; modelled as all nodes on the row
$y = 2200$\unit{mm}, imposed horizontal displacement \\
Control node   & model convention & $(x,y) = (0, 2200)$\unit{mm} \\
Drift ratio    & model convention & control displacement divided by 2200\unit{mm} \\
Base shear     & model convention & $-\sum$ horizontal reactions over the pedestal soffit nodes \\
\bottomrule
\end{tabularx}

Gravity is applied first in 10 \texttt{LoadControl} increments and then held constant
(\texttt{loadConst}) for the whole cyclic stage.

\section{Loading protocol}

The specimen's protocol, from paper Fig.~9 and Sec.~2.2 (after ACI 374.2R-13): three cycles at each
level up to the yield displacement and two at every level after it. Amplitudes are measured at the
actuator level, $y = 2200$\unit{mm}, and the drift levels 0.3--4.0\% and the 2.2 / 4.4\unit{mm}
amplitudes below them are read directly from Fig.~9, where 4.4\unit{mm} is annotated $\Delta_y$ and
2.2\unit{mm} is $0.5\Delta_y$.

The model reproduces the displacement-controlled phase only. The test additionally began with
force-controlled cycles raised to the cracking load of 45.4\unit{kN} ($M_{cr}$, Sec.~2.2); those are
omitted, since a displacement-driven analysis cannot impose them and they precede any measurable
inelasticity.

\begin{center}\small
\begin{tabular}{@{}lrrrrrrrrrrr@{}}
\toprule
Drift ratio (\%) & 0.10 & 0.20 & 0.30 & 0.50 & 0.75 & 1.00 & 1.50 & 2.00 & 2.50 & 3.00 & 4.00 \\
Amplitude (mm)   & 2.2 & 4.4 & 6.6 & 11.0 & 16.5 & 22.0 & 33.0 & 44.0 & 55.0 & 66.0 & 88.0 \\
Cycles           & 3 & 3 & 2 & 2 & 2 & 2 & 2 & 2 & 2 & 2 & 2 \\
\bottomrule
\end{tabular}\end{center}

Totals: 11 levels, 24 cycles, 49 reversals, 2818\unit{mm} of accumulated displacement path.

\section{Solution procedure}

The solution procedure has no counterpart in the paper --- the specimen was loaded quasi-statically
by hydraulic actuator --- so every setting below is a numerical choice made to trace that same
history. The cyclic response is traced by dynamic relaxation: a quasi-static transient solve in
which the displacement history is imposed through single-point constraints and the drive is slow
enough that inertial and damping contributions to the recorded reactions are negligible.

\begin{tabularx}{\textwidth}{@{}p{50mm}p{38mm}>{\raggedright\arraybackslash}X@{}}
\toprule
\textbf{Setting} & \textbf{Value} & \textbf{Note} \\
\midrule
Analysis type         & \texttt{Transient} & \\
Displacement history  & \texttt{Path} time series, imposed via \texttt{sp} constraints & one continuous solve across all reversals \\
Drive rate            & 7.6\unit{mm/s}, constant & \\
Fundamental period $T_1$ & 0.0240\unit{s} & from \texttt{eigen}, gravity-held tangent \\
Time step $\Delta t$  & $T_1/30 = 8.0\times10^{-4}$\unit{s} & \\
Total steps           & 463\,534 & \\
Integrator            & \texttt{HHT}, $\alpha = 0.7$ & numerical damping of high-frequency content \\
Damping               & Rayleigh, $\zeta = 0.8$ & $\alpha_M = \zeta\omega_1 = 209.4$, $\beta_{K,\text{init}} = \zeta/\omega_1 = 3.056\times10^{-3}$ \\
Damping stiffness     & \textbf{initial} stiffness ($\beta_{K,\text{init}}$), not the current tangent & \\
System                & \texttt{BandGeneral} & \\
Numberer              & \texttt{RCM} & \\
Constraints           & \texttt{Transformation} & \\
Convergence test      & \texttt{NormDispIncr}, tol $10^{-5}$, 50 iterations & \\
Primary algorithm     & \texttt{Newton} & \\
Fallback algorithms   & \texttt{KrylovNewton}, then \texttt{NewtonLineSearch} & retried at $\Delta t/10$ over 10 sub-steps \\
\bottomrule
\end{tabularx}

\texttt{wipeAnalysis} is called after the gravity stage and before the transient integrator is set;
the solver (constraints, numberer, system) is then re-declared. Without this the gravity stage's
static analysis remains active and the transient integrator is not accepted.

The drive rate is specified directly rather than derived from a number of fundamental periods to
the target amplitude, so that the speed is independent of the protocol's amplitude. At
7.6\unit{mm/s} the recorded peak base shear reproduces that of the static displacement-controlled
solver to within 1\% at 0.10\% drift.

\section{Idealizations}\label{sec:ideal}

\begin{tabularx}{\textwidth}{@{}p{46mm}>{\raggedright\arraybackslash}X@{}}
\toprule
\textbf{Item} & \textbf{Treatment} \\
\midrule
Pedestal thickness & 800\unit{mm} against the wall's 200\unit{mm}. A plane model carries one
thickness, so the pedestal's extra width is folded into an equivalent modulus,
$E \times 800/200 = 93\,600$\unit{MPa}. \\
Pedestal material & linear elastic throughout. \\
Boundary bar group & 4$\phi$14 lumped at the group centroid, preserving the first moment of the
group about the section. \\
Reinforcement curtains & two curtains collapsed onto one in-plane line at double area. \\
Hoop geometry & only in-plane legs are representable in a 2D in-plane model; legs running through
the thickness are out of plane. \\
Bond & perfect; reinforcement shares lattice nodes with the concrete. \\
Bar buckling and fracture & not represented by \texttt{Steel02}. \\
Concrete compression & \texttt{Concrete02} crushing backbone. A selectable mode replaces it with an
unbounded elastic branch, matching the assumption of the cited calibration method, in which
compression is not calibrated. \\
\bottomrule
\end{tabularx}

\section{Outputs}

Recorded at every converged step: control-node displacement and base shear. Written to
\path{examples/output/wall/}:

\small
\begin{tabularx}{\textwidth}{@{}p{62mm}>{\raggedright\arraybackslash}X@{}}
\toprule
\textbf{File} & \textbf{Content} \\
\midrule
\texttt{wall\_cyclic\_dynamic.png} & hysteresis loops, measured response, overlay \\
\texttt{wall\_cyclic\_dynamic\_backbone.png} & loop envelope against the measured backbone \\
\texttt{wall\_cyclic\_dynamic\_vs\_test.png} & loops and envelopes against the \emph{digitized} test response \\
\texttt{wall\_cyclic\_dynamic\_data.json} & raw drift and base-shear arrays \\
\texttt{wall\_drawing.png} & lattice model elevation and interface detail \\
\bottomrule
\end{tabularx}
\normalsize

The digitized test response is kept with the study rather than the output, at
\path{examples/wall/data/swncff_fig14b.npz}, since it is an input to the comparison and does not
change when the analysis is re-run.

Measured reference values for this specimen: peak lateral load $+212.6\,/\,-209.4$\unit{kN} at
$\pm1.00\%$ drift; first bar yield at 0.30\% drift; idealized yield drift 0.35\%; ultimate drift
1.50\%; displacement ductility 4.3; strength degradation 35\% at 4.0\% drift.

\begin{center}
\includegraphics[width=0.96\textwidth]{wall_drawing.png}
\captionof{figure}{Lattice model of Kutay's wall: full elevation and wall--pedestal interface.}
\end{center}

\section{Results}

The analysis completed all 49 reversals of the 4.0\% protocol, 463\,534 steps in 7.1\unit{h},
converged throughout.

\subsection{Digitization of the measured response}

The measured hysteresis is taken from Figure 14b of the source paper (specimen SW-NC-FF). The page
is rendered at 400\unit{dpi}, the plot frame is located from its own axis rules, and the test curve
is separated from the damage markers by colour, giving 64\,414 points in (drift ratio, lateral
load). Both the model and the digitized response use the same shear span, 2200\unit{mm}, so the two
overlay without rescaling.

The extraction is self-checking: the digitized extremes are $+210.6\,/\,-211.8$\unit{kN} against
the reported $+212.6\,/\,-209.4$\unit{kN}, a 1.0\% and 1.2\% difference, against a plotted line
thickness worth roughly 3\unit{kN}. What is recovered is a point \emph{cloud}, not an ordered path:
overlapping loops cannot be re-sequenced from pixels. Per-cycle quantities --- dissipated energy,
cycle-to-cycle degradation, unloading stiffness --- are therefore not available from this source,
and the cloud is plotted with markers rather than joined by lines.

A backbone is recovered as a \emph{loop-tip envelope}: at each protocol drift level $d$ the extreme
load is taken over the slice $[0.90d,\ 1.05d]$. The slice reaches 10\% inside the extreme because
the measured load peaks before the peak displacement of the same cycle. Wider inner reaches admit
the reloading branches of larger loops: at 15\% the 3.0\% level reads 174\unit{kN} instead of
152\unit{kN}, taken from the 4\% cycle. A per-drift-bin maximum reads 180\unit{kN} at 0.3\% drift,
also from larger cycles.

\subsection{Peak strength}

\begin{tabularx}{\textwidth}{@{}p{36mm}p{24mm}p{24mm}p{26mm}>{\raggedright\arraybackslash}X@{}}
\toprule
\textbf{Quantity} & \textbf{Push} & \textbf{Pull} & \textbf{At drift} & \textbf{Source} \\
\midrule
Model peak base shear & $+227.2$\unit{kN} & $-221.5$\unit{kN} & $+0.36$ / $-0.30\%$ & this analysis \\
Test, digitized       & $+209.9$\unit{kN} & $-211.8$\unit{kN} & $+1.00$ / $-1.50\%$ & Fig.~14b, digitized \\
Test, as reported     & $+212.6$\unit{kN} & $-209.4$\unit{kN} & $\pm1.00\%$ & paper Sec.~3 \\
Test $V_u$            & \multicolumn{2}{l}{211\unit{kN} (average)} & --- & paper Table~3 \\
Model / test          & 1.08 & 1.05 & --- & against the digitized peaks \\
\bottomrule
\end{tabularx}

\subsection{Envelope, level by level}

Model envelope against the digitized loop-tip envelope, at each protocol drift level (kN):

\begin{center}\small
\begin{tabular}{@{}lrrrrrrrrr@{}}
\toprule
Drift ratio (\%) & 0.30 & 0.50 & 0.75 & 1.00 & 1.50 & 2.00 & 2.50 & 3.00 & 4.00 \\
\midrule
Model, push      & 218.3 & 208.7 & 207.1 & 202.2 & 184.7 & 187.5 & 185.7 & 179.1 & 189.6 \\
Test, push       & 165.5 & 203.7 & 204.3 & 209.9 & 207.4 & 202.4 & 174.9 & 152.4 & 141.7 \\
Model / test     & 1.32 & 1.02 & 1.01 & 0.96 & 0.89 & 0.93 & 1.06 & 1.18 & 1.34 \\
\midrule
Model, pull      & $-212.9$ & $-211.0$ & $-209.5$ & $-205.3$ & $-200.2$ & $-188.5$ & $-183.7$ & $-178.7$ & $-182.7$ \\
Test, pull       & $-163.0$ & $-199.9$ & $-201.8$ & $-209.9$ & $-211.8$ & $-196.8$ & $-171.8$ & $-145.5$ & $-134.2$ \\
Model / test     & 1.31 & 1.06 & 1.04 & 0.98 & 0.95 & 0.96 & 1.07 & 1.23 & 1.36 \\
\bottomrule
\end{tabular}\end{center}

Between 0.50\% and 2.00\% drift the envelopes agree to within 11\%, and over 0.50--1.00\% drift to
within 6\%.

\subsection{Measured quantities not reproduced by the model}

Each entry states the reported value, the modelled value, and the model feature that carries the
mechanism. Section~\ref{sec:ideal} lists the corresponding idealizations.

\begin{tabularx}{\textwidth}{@{}p{27mm}p{37mm}p{25mm}>{\raggedright\arraybackslash}X@{}}
\toprule
\textbf{Quantity} & \textbf{Reported} & \textbf{Model} & \textbf{Mechanism in the model} \\
\midrule
Drift at peak load & 1.00\% (Sec.~3) & 0.36\% & bond is perfect; bar slip is not represented \\
Strength at 4.0\% drift & 35\% degradation, $\approx138$\unit{kN} (Sec.~5) & 189.6\unit{kN}, no degradation &
\texttt{Steel02} has no bar buckling; the reported degradation is bar buckling with core crushing (Sec.~3) \\
Ultimate drift $\delta_u$ & 1.50\% (Table~3) & not reached by 4.0\% & no strength-loss criterion is modelled \\
Displacement ductility $\mu$ & 4.3 (Sec.~3) & not evaluated & follows from $\delta_u$ \\
Deformation split at 2.0\% drift & 74\% rocking, 20\% flexure, 4\% shear, 2\% other (Sec.~4) &
flexure and shear only & rocking arises from debonding of plain bars \\
Loop shape & pinched, self-centering (Sec.~3) & wider loops &
same mechanism as above \\
Per-cycle dissipated energy & --- & --- & not comparable: the digitized source is an unordered point cloud \\
\bottomrule
\end{tabularx}

The lateral drive is displacement-controlled and the protocol closes at zero displacement, so the
end-of-run displacement is imposed rather than computed.

\begin{center}
\includegraphics[width=\textwidth]{wall_cyclic_dynamic_vs_test.png}
\captionof{figure}{Lattice response against the digitized SW-NC-FF test hysteresis. Left: loops
overlaid, the test drawn as a point cloud. Right: the loop envelopes.}
\end{center}

\end{document}
